\textit{Air Balloon Survival} is a conversational negotiation game which aims to evaluate advanced reasoning and the degree of interactive collaboration capabilities. Two players are travelling in a sinking hot air balloon and must agree to keep a set of items such that they do not exceed the balloon's maximum weight capacity. Each player is presented with the list of available items and their respective weight. Additionally, both players receive unique preference values for each item which remain hidden from the other player. Players must then negotiate and make a deal which maximizes the utility score, i.e. the sum of preference values of items while respecting the weight constraint. During negotiation, players must respect the rules of the game which include a clear framework governing the syntactic structure of their messages and rules governing which actions they are allowed to take to advance in the game. The game ends once a deal has been made. A deal is defined as an explicit agreement by one player with a proposal their interlocutor has made. A proposal can be any subset of the set of items in the game. A simplified example of an episode can be seen in Figure \ref{fig:air_balloon_example}. \textit{Air Balloon Survival} probes the players abilities to reason and negotiate on two levels. On (1) the level of each individual player and (2) on a collective level. (1) On an individual level, each player is effectively given an instance of a version of the 0/1 Knapsack Problem. Hence, on this individual level our game probes a players ability to perform constraint-based optimization which includes carrying out arithmetic operations and performing combinatorial search over all possible deals. (2) On the collective level, we probe the players’ practical rationality, i.e. their ability to maximize personal utility by successfully negotiating an agreement with their counterpart who may have conflicting goals. Given that players must align by inferring their counterparts’ implicit preferences and intentions from observed responses, the game also connects to the broader topic of Theory of Mind (ToM) and LLMs.
\begin{figure}[ht!]
    \centering
    % \newcounter{utterance}
{\footnotesize  \setcounter{utterance}{1}
\setlength{\tabcolsep}{0pt}
\begin{supertabular}{c@{$\;$}|p{.15\linewidth}@{}p{.15\linewidth}p{.15\linewidth}p{.15\linewidth}p{.15\linewidth}p{.15\linewidth}}   
\# & $\;$A & \multicolumn{4}{c}{Game Master} & $\;\:$B\\
    \hline

\theutterance \stepcounter{utterance}  
    & & \multicolumn{4}{p{0.6\linewidth}}{
        \cellcolor[rgb]{0.9,0.9,0.9}{
            \makecell[{{p{\linewidth}}}]{
                \texttt{\tiny{[A$\langle$GM]}}
                <GAME DESCRIPTION>\\
                \texttt{-Items and weights: Ball(2), Book(3), Magazine(1), Rope(4), Lamp(5)} \\
                \texttt{-Max Weight: 7} \\
                \texttt{-Your Preferences: Ball=1, Book=2, Magazine=2, Rope=2, Lamp=1}
            }
        }
    }
    & \\
    
\theutterance \stepcounter{utterance}  
    & \multicolumn{4}{p{0.6\linewidth}}{
        \cellcolor[rgb]{1,0.9,0.9}{
            \makecell[{{p{\linewidth}}}]{
                \texttt{\tiny{[A$\rangle$GM]}}
                \texttt{PROPOSAL: Magazine, Ball, Book} \\
                \texttt{ARGUMENT: These items balance weight and value well for me. What do you think?}
            }
        }
    }
    & & \\

\theutterance \stepcounter{utterance}  
    & & \multicolumn{4}{p{0.6\linewidth}}{
        \cellcolor[rgb]{0.9,0.9,0.9}{
            \makecell[{{p{\linewidth}}}]{
                \texttt{\tiny{[GM$\rangle$B]}}
                <GAME DESCRIPTION>\\
                \texttt{<Items, weights, weight limit>}\\
                \texttt{-Your Preferences: Ball=1, Book=1, Magazine=2, Rope=4, Lamp=5} \\
                \centering{<PLAYER A Message>}
            }
        }
    }
    & \\

\theutterance \stepcounter{utterance}  
    & & & \multicolumn{4}{p{0.6\linewidth}}{
        \cellcolor[rgb]{0.9,1,0.9}{
            \makecell[{{p{\linewidth}}}]{
                \texttt{\tiny{[GM$\langle$B]}}
                \texttt{PROPOSAL: Rope, Ball, Magazine}\\
                \texttt{ARGUMENT: I value Rope highly, but not Magazine. How about we swap them?}
            }
        }
    } \\

\theutterance \stepcounter{utterance}  
    & & \multicolumn{4}{p{0.6\linewidth}}{
        \cellcolor[rgb]{0.9,0.9,0.9}{
            \makecell[{{p{\linewidth}}}]{
                \texttt{\tiny{[A$\langle$GM]}}
                \centering{<PLAYER B Message>}
            }
        }
    }
    & \\

\theutterance \stepcounter{utterance}  
    & \multicolumn{4}{p{0.6\linewidth}}{
        \cellcolor[rgb]{1,0.9,0.9}{
            \makecell[{{p{\linewidth}}}]{
                \texttt{\tiny{[A$\rangle$GM]}}
                \texttt{AGREE: Rope, Ball, Magazine} \\
                \texttt{ARGUMENT: I am indifferent w.r.t. Rope and Magazine. Fair deal.} \\
            }
        }
    }
    & & \\
    
\end{supertabular}
}
    \caption{A simplified dialogue example from the \textit{Air Balloon Survival} game. In this game, two players must negotiate and argue for their preferred set of items. To reach the goal state, one player must explicitly agree to a proposal made by the other. Messages from player A and player B are shown in red and green, respectively. Grey messages come from the game master, who provides the initial prompts and relays messages between the two players. \tk{I'd say add some made-up list of items and their weights from each player's view so that we see what is there. and they are trying to negotiate for/about etc. Maybe Player 2 could accept some items but negotiate over others based on personal preferences (weights) etc. so that we see how personal preferences play a role here.}}
    \label{fig:air_balloon_example}
\end{figure}
\paragraph{Game Mechanics}
Both players receive initial prompts that may differ in the preference values assigned to items. They are instructed to use specified formats when making proposals and negotiating, and to follow rules such as only being allowed to accept proposals made by the other player (cf. Fig.~\ref{fig:hot_air_balloon_init_prompt} in Appendix~\ref{sec:hot_air_balloon_appendix_prompt_templates}). 
There is no limit on the number of turns, but a game is aborted if no progress is made for eight consecutive turns. Each player may commit up to two mistakes, either by violating game rules or by violating the required syntax. In such cases, they are reprompted with a tailored message (cf. Figs.~\ref{fig:hot_air_balloon_parse_error_prompts}, \ref{fig:hot_air_balloon_invalid_action_prompts}). 
Messages are relayed through the game master. In some of our game instances we allow players to use a \texttt{STRATEGIC REASONING} tag the contents of which will not be forwarded to the other player by the game master which allows a player to `think out loud' about, for example, the counterpart's preferences or their own explicit preference values. When used with a reasoning model, the tag can be compared to the model’s internal reasoning traces. 
A game ends successfully once a player accepts a proposal made by the counterpart.
\paragraph{Instance Generation}
For instance generation, we draw either 15 or 35 items (depending on the experiment) from a randomly generated list constructed by concatenating a capital letter with a two-digit number (e.g., \texttt{A42}, \texttt{C07}). For each item $i$ we draw a weight $w_i$ from some discrete uniform distribution. We used $U\{1,1000\}$ in all instances. The air balloon's capacity is defined as a fraction of the combined weight of all items, i.e., $W = \alpha \sum_{i=1}^{n} w_i$, where $\alpha \in (0,1)$ is a scaling parameter controlling the tightness of the weight constraint. In all of our instances we set $\alpha = 0.5$, so the capacity of the air balloon equals half of the sum of the weights. For each player $p \in P$ we generate an individual valuation function assigning a preference value $v_p(i)$ to each item $i$. We control two dimensions of correlation when generating these valuations: (i) the correlation between weights and preferences, and (ii) the correlation between players’ preferences. 
For (i) we distinguish two cases. Firstly, there is the uncorrelated case, where preference values are drawn independently of the weights, $v_p(i) \sim U\{1,R\}, \ \forall i \in \{1,\dots,n\}$, where $R$ denotes the maximum weight among all items. Secondly, there is the strongly correlated case, where the preference values are defined directly from the weights, $v_p(i) = w_i + q, \ \forall i \in \{1,\dots,n\}$, with $q = R/10$ controlling the correlation offset. For (ii) we distinguish three cases. In the equal case, both players share identical valuation functions, $v_1(i) = v_2(i)$ for all $i$, thus giving them the same goals. In the independent case, the valuations $v_1$ and $v_2$ are generated separately. In the inverse case, the rankings of the items are reversed across players, so that the highest-valued item of player 1 receives the lowest value for player 2, and so forth, giving them opposing goals. 
Increasing the number of items and controlling the correlation between weights and preferences relates to (1), the hardness of the knapsack problem and the reasoning required on an individual player’s level. The correlation between players’ preferences relates to (2), their ability to negotiate, take appropriate actions and reach an agreement when they have common or opposing goals.
\paragraph{Evaluation Metrics}
Each player $p \in P$ receives an individual score based on the utility of the final deal $D$, i.e. the sum of the preference values of items in $D$. This sum gets normalized by the value of the optimal solution to their instance of the 0/1 Knapsack Problem. The overall game score is defined as the harmonic mean of the two players’ normalized scores, further normalized by the maximum harmonic mean achievable for an instance. All scores are scaled to the range $[0,100]$. If the total weight of the deal exceeds the capacity $W$, the score is set to $0$. We chose the harmonic mean as our main metric 
as it rewards collectively balanced outcomes by favouring deals where both players achieve high scores and penalizing those where one player benefits disproportionately. 
The choice for individual scores is straightforward, as they simply measure each player’s outcome relative to the optimal solution of their own knapsack problem.
For a final deal $D \subseteq \{1,\dots,n\}$, the normalized score of player $p \in P$ is  
$$
f_p(D) = 100 \cdot \frac{\sum_{i \in D} v_p(i)}{\text{OPT}_p},
$$  
where $\text{OPT}_p$ is the value of the optimal solution to player $p$'s 0/1 Knapsack Problem.  
The harmonic mean of the players' normalized scores is defined as  
$$
f_{\text{harm}}(D) =
\begin{cases}
\frac{2 f_{1}(D) f_{2}(D)}{f_{1}(D) + f_{2}(D)} & \text{if } \sum_{i \in D} w_i \leq W, \\
0 & \text{otherwise}.
\end{cases}
$$  
The final game score is the normalized harmonic mean, 
$$
\text{score}(D) = 100 \cdot \frac{f_{\text{harm}}(D)}{\text{OPT}^*},
$$  
where $\text{OPT}^*$ denotes the maximum harmonic mean achievable in the given instance.
\tk{Game desc etc.}